\chapter{Introduction à la Gestion des Incidents}
\label{chap:introincidents}

Malgré toutes les mesures de sécurité (Prévention) et les vérifications (Audit), aucune organisation n'est infaillible. La question n'est plus "Est-ce que nous allons subir une attaque ?", mais "Quand allons-nous subir une attaque et comment réagir ?". Ce chapitre pose les bases de la réponse opérationnelle.

\section{Présentation : La Nécessité de Réagir}

\subsection{De la Sécurité Préventive à la Sécurité Curative}

Jusqu'ici, nous avons travaillé sur la prévention (analyse de risque, mesures de sécurité, audits). Cependant, le "Risque Résiduel" (Module 1) implique que des incidents surviendront.

La capacité à détecter et réagir efficacement est ce qui différencie une entreprise résiliente d'une entreprise qui disparaît suite à une cyberattaque.

\subsection{La Pyramide de l'Événement}

Il est crucial de distinguer les concepts pour ne pas noyer les équipes sous de fausses alertes.

\begin{figure}[htbp]
\centering
\begin{tikzpicture}[scale=0.8]
    % Triangle inversé
    \node[draw, isosceles triangle, isosceles triangle apex angle=30, 
          minimum height=6cm, minimum width=8cm, 
          shape border rotate=90, draw=black, fill=gray!10] at (0,0) {};
    
    % Labels
    \node at (0, 2.5) {\textbf{Crise}};
    \node at (0, 1) {\textbf{Incident}};
    \node at (0, -1) {\textbf{Événement}};
    
    % Annotations
    \node[right] at (5, 2.5) {Impact Majeur / Existence en jeu};
    \node[right] at (5, 1) {Impact Réel / Mesures requises};
    \node[right] at (5, -1) {Observation brute / Neutre};
\end{tikzpicture}
\caption{Pyramide Événement, Incident, Crise}
\end{figure}

\begin{boxdef}
\begin{itemize}
    \item \textbf{Événement (Event) :} Un changement d’état observé dans un système. La majorité des événements sont normaux (ex: un utilisateur se connecte, un email arrive).
    \item \textbf{Incident :} Un événement qui compromet réellement ou potentiellement la sécurité (Confidentialité, Intégrité, Disponibilité). Il nécessite une réponse.
    \item \textbf{Crise :} Un incident dont l'impact dépasse la capacité de traitement opérationnel de l'entreprise et menace sa survie ou sa réputation.
\end{itemize}
\end{boxdef}

\begin{boxlizbiz}
\textbf{Exemples chez LiZbiZ :}
\begin{itemize}
    \item \textbf{Événement :} Un employé saisit son mot de passe. (Normal).
    \item \textbf{Incident :} Un employé clique sur un lien de phishing et ses identifiants sont volés. (Nécessite une action immédiate : changement de mdp, analyse).
    \item \textbf{Crise :} Le Ransomware se propage, l'ERP est inutilisable, la presse s'en empare. (Nécessite une cellule de crise).
\end{itemize}
\end{boxlizbiz}

% ------------------------------------------------------------------------------
\section{Concepts Clés}
% ------------------------------------------------------------------------------

\subsection{Cycle de Vie de l'Incident}

La gestion d'incident ne s'improvise pas. Elle suit un cycle standard (souvent inspiré du NIST SP 800-61) :

\begin{enumerate}
    \item \textbf{Détection et Analyse :} Repérer l'incident (Logs, Alerte utilisateur, Antivirus).
    \item \textbf{Confinement :} Empêcher la propagation (Isolation réseau, Désactivation compte).
    \item \textbf{Éradication :} Supprimer la cause (Supprimer malware, Patch faille).
    \item \textbf{Rétablissement :} Remettre le système en marche (Restauration sauvegarde).
    \item \textbf{Activités Post-Incident :} Retour d'expérience (RETEX), amélioration des défenses.
\end{enumerate}

\subsection{Impact Business vs Impact Technique}

Un bon gestionnaire d'incident traduit l'impact technique en impact métier.

\begin{table}[htbp]
\centering
\begin{tabular}{p{4cm}p{4cm}}
    \toprule
    \textbf{Impact Technique} & \textbf{Impact Business (LiZbiZ)} \\
    \midrule
    "Le serveur Web est down" & "Les clients ne peuvent plus commander (Perte de 5k€/heure)". \\
    "Fuite de la BDD Clients" & "Risque d'amende CNIL, perte de confiance, baisse du cours de bourse". \\
    \bottomrule
\end{tabular}
\caption{Traduction Technique vers Business}
\end{table}

\subsection{Obligation de Reporting}

Depuis le RGPD et la directive NIS, certaines infractions doivent être notifiées.

\begin{boxdef}
\textbf{Notification de violation de données (RGPD Art. 33) :} En cas de violation de données personnelles, l'entreprise dispose de \textbf{72 heures} pour notifier l'autorité de contrôle (CNIL en France), sous peine de sanction financière lourde.
\end{boxdef}

% ------------------------------------------------------------------------------
\section{Jargon de la Gestion d'Incident}
% ------------------------------------------------------------------------------

\subsection{Alerte et Ticket}

\begin{itemize}
    \item \textbf{Alerte :} Signal automatique généré par un outil (SIEM, Antivirus) indiquant une anomalie potentielle.
    \item \textbf{Ticket :} Enregistrement formel de l'incident dans un système de ticketing (Jira, GLPI, ServiceNow) pour le suivi.
\end{itemize}

\subsection{Faux Positif (False Positive)}

\begin{boxdef}
\textbf{Faux Positif :} Une alerte qui signale une menace alors qu'il n'y en a pas. C'est le fléau des équipes SOC (Security Operations Center). Un taux de faux positifs trop élevé épuise les analystes qui finissent par ignorer les vraies alertes.
\end{boxdef}

\begin{boxlizbiz}
\textbf{LiZbiZ :} L'antivirus signale un "comportement suspect" sur l'outil de caisse. Après vérification (Audit), on s'aperçoit que c'est une mise à jour logicielle légitime mal signée. C'est un Faux Positif.
\end{boxlizbiz}

% ------------------------------------------------------------------------------
\section{Cas LiZbiZ : Périmètre et Classification}
% ------------------------------------------------------------------------------

\subsection{Définition du Périmètre}

Pour LiZbiZ, tout événement affectant la disponibilité du site e-commerce ou l'intégrité des données clients est considéré comme un incident de sécurité.

\subsection{Exercice : Classification des Incidents}

Proposez une classification (Critique, Haute, Basse) pour les incidents suivants chez LiZbiZ :

\begin{enumerate}
    \item Un employé a perdu sa badge d'accès au bâtiment.
    \item Le serveur de base de données est inaccessible.
    \item Une publicité pop-up suspecte apparaît sur un site web tiers visité par un employé.
    \item Un email de phishing demandant les identifiants ERP est envoyé à tout le service comptabilité.
\end{enumerate}

\begin{boxlizbiz}
\textbf{Proposition de solution :}
\begin{itemize}
    \item Badge perdu : \textbf{Basse} (Mesure : Désactiver le badge).
    \item BDD inaccessible : \textbf{Critique} (Arrêt business).
    \item Pop-up tiers : \textbf{Négligeable} (Pas d'impact direct sur le SI LiZbiZ).
    \item Phishing comptabilité : \textbf{Haute} (Risque de fraude financière, nécessite une sensibilisation rapide et vérification des clics).
\end{itemize}
\end{boxlizbiz}

% ==============================================================================
% Fichier : 04_module_incident.tex (Extrait Chapitre 2)
% Description : Le Processus de Gestion d'Incident (ISO 27035)
% ==============================================================================